\documentclass[12pt]{article}
\usepackage[utf8]{inputenc}
\usepackage[T1]{fontenc}
\usepackage[margin=1in]{geometry}
\usepackage{setspace}
\usepackage{xcolor}
\usepackage[colorlinks=true, urlcolor=blue!60!black, linkcolor=blue!60!black]{hyperref}
\setlength{\parindent}{0em}
\setlength{\parskip}{0.6em}
\renewcommand{\baselinestretch}{1.0}
\begin{document}
\begin{center}
{\LARGE \textit{Christian Legal Society v.\ Martinez}, 561 U.S.\ 661 (2010)\par}
\vspace{0.5em}
{\large Case Brief --- Two-Page Edition\par}
\end{center}

\section{Memory Jogger}
A public law school may condition student-group recognition on an ``all-comers'' membership policy, a reasonable, viewpoint-neutral condition on a limited public forum (\href{https://supreme.justia.com/cases/federal/us/561/661/}{561 U.S. 661}).

\section{Detailed Case Facts}
Hastings' Registered Student Organization program offered funding from student fees, official communication channels, and use of facilities and the school's name. The parties stipulated that Hastings required every group to accept all comers as members and leaders. The Christian Legal Society chapter required members and officers to sign a Statement of Faith and excluded those engaged in ``unrepentant homosexual conduct.'' Hastings denied it recognition but let it use rooms and general bulletin boards.

\section{Procedural History}
CLS sued under \S~1983. The district court granted summary judgment to Hastings; the Ninth Circuit summarily affirmed; the Supreme Court affirmed and remanded the pretext claim.

\section{Judicial Votes}
5--4. Ginsburg wrote for Stevens, Kennedy, Breyer, and Sotomayor; Stevens and Kennedy concurred; Alito dissented, joined by Roberts, Scalia, and Thomas.

\section{Holding}
Speech and expressive-association challenges to a limited forum's membership conditions are reviewed together under forum doctrine, not strict scrutiny. A genuine all-comers policy is reasonable and viewpoint neutral, and \textit{Smith} bars a free-exercise exemption.

\section{Analysis of Opinions}
Holding CLS to its stipulation, Ginsburg reasoned that Hastings offered a subsidy rather than a prohibition, so \textit{Jaycees} and \textit{Dale} did not apply. The policy opens fee-funded leadership to all, avoids policing belief-versus-status lines, encourages tolerance, and tracks state law; it is ``textbook viewpoint neutral'' despite uneven effects. Kennedy warned that infiltration or pretext would present ``a substantial case on the merits.'' Alito saw pretext and viewpoint discrimination, invoked \textit{Healy v.\ James}, and warned of hostile takeovers.

\section{Examples: Future Applications}
\textbf{Same-Side:} a uniformly enforced all-comers rule denying funds to a club with an ideological officer pledge is upheld. \textbf{Opposite-Side:} deregistering a Christian group while letting secular groups require mission alignment is viewpoint discrimination. \textbf{Fence-Sitter:} opponents use an all-comers rule to take over a small religious club.

\section{Critique}
Critics say the stipulation produced an artificial case and that all-comers rules defeat a forum meant to foster diverse groups; originalists see the right of assembly subordinated. Progressives largely defend the result as civic equality, though some worry about their own affinity groups.

\section{Key Quotations}
\begin{quote}
``The First Amendment shields CLS against state prohibition of the organization's expressive activity, however exclusionary that activity may be. But CLS enjoys no constitutional right to state subvention of its selectivity.''
\end{quote}
\noindent --- Justice Ginsburg, Opinion of the Court, 561 U.S. at 669.

\end{document}
